\documentclass[FGBook.tex]{subfiles}

\begin{document}

\proza{Harry Potter a Kámen mudrců}{JOANNE ROWLINGOVÁ}

\noindent\begin{wrapfigure}{r}{0.35\textwidth}\footnotesize
\subwection{Ukázka z díla:}
\setlength{\parindent}{3pt}
„Jsem čaroděj,“ řekl Harry.\\
„Kouzelník, samozřejmě,“ přikývl Hagrid a posadil se na pohovku, která pod jeho vahou zapraskala. „A to ještě pořádný, jak slyším. S tou jizvou a tak. Ne, ty jsi slavný, Harry.“
\end{wrapfigure}

\wection{LITERÁRNÍ TEORIE}

\subwection{Literární druh a žánr:}
\noindent epická próza, fantasy román (dobrodružný školní příběh s prvky coming-of-age)

\subwection{Literární směr:}
\noindent Současná dětská a young adult literatura s prvky moderního fantasy. Inspirace britskou tradicí školních příběhů, pohádek a mytologie.

\subwection{Jazyk a slovní zásoba:}
\noindent Přístupný, živý a imaginativní jazyk s bohatou slovní zásobou vymyšlených kouzelnických termínů (bradavický argot, jména jako Bradavice, Nebelvír). Mnoho popisů, humoru a dialogů. Spisovný jazyk s prvky hovorovosti u postav (Hagrid).

\subwection{Figury:}
\noindent Dialogy, opakování motivů (jizva, jméno „Ten, jehož jméno nesmíme vyslovit“), ironie a humor.

\subwection{Tropy:}
\noindent Symbolismus (\texthl{jizva = osud a přežití}), metafora (Kámen mudrců = nesmrtelnost a chamtivost), alegorie (boj dobra proti zlu), personifikace kouzelných předmětů.

\subwection{Stylistická charakteristika textu:}
\noindent Chronologické vyprávění v er-formě se vševědoucím vypravěčem. Rychlý spád, střídání napětí, humoru a dojemných momentů. Detailní popisy kouzelného světa kontrastují s šedivým mudlovským životem.

\subwection{Postavy:}
\noindent 
\textbfhl{HARRY POTTER --} hlavní hrdina, sirotek s jizvou ve tvaru blesku, statečný, loajální, skromný. Přechází z týraného chlapce k hrdinovi.\\
\textbfhl{RON WEASLEY --} věrný přítel z chudé kouzelnické rodiny, odvážný, vtipný, někdy nejistý.\\
\textbfhl{HERMIONA GRANGEROVÁ --} chytrá, pilná, z mudlovské rodiny. Nejlepší studentka, často zachraňuje situaci znalostmi.\\
\textbfhl{ALBUS BRUMBÁL --} moudrý a laskavý ředitel Bradavic, ochránce Harryho, symbol dobra.\\
\textbfhl(PROFESOR QUIRRELL --} učitel obrany proti černé magii, slabý a koktavý, ve skutečnosti posedlý Voldemortem.

\subwection{Děj:}
\noindent Harry Potter žije u Dursleyových, kteří ho nenávidí. Na jedenácté narozeniny se dozví, že je čaroděj a je přijat do Bradavic. Seznámí se s Ronem a Hermionou, stane se členem Nebelvíru. Ve škole objeví, že někdo chce ukrást Kámen mudrců – artefakt umožňující nesmrtelnost. Společně s přáteli překoná řadu pastí (tříhlavý pes, ďábelské pasti atd.) a zjistí, že za tím stojí profesor Quirrell posedlý lordem Voldemortem. Harry kámen zachrání díky lásce své matky, která ho chrání před zlem.

\subwection{Kompozice, prostor a čas:}
\noindent Chronologický děj s retrospektivami do Harryho dětství. Prostor: mudlovský svět (Privet Drive) × kouzelný svět (Diagon Alley, Bradavice, nástupiště 9¾). Čas: konec 20. století (Harrymu je 11 let).

\subwection{Význam sdělení:}
\noindent Boj dobra proti zlu, síla přátelství, odvahy a lásky. Důležitost volby („naši volby ukazují, kým doopravdy jsme“). Kontrast mezi šedivým mudlovským světem a barevným kouzelným světem symbolizuje objevování vlastní identity a překonávání osudu.

\vspace*{2em}\\
\wection{LITERÁRNÍ HISTORIE}

\subwection{Politická situace:}
\noindent Kniha vyšla v roce 1997 – období relativního míru v Británii, ale s rostoucím zájmem o fantasy a escapismus.

\subwection{Autor {\ssmall -- život autora:}}
\noindent Joanne Kathleen Rowlingová (*1965). Britská spisovatelka. Nápad na Harryho dostala v roce 1990 ve vlaku. Psala v kavárnách jako samoživitelka. První díl byl odmítnut mnoha nakladateli, pak uspěl u Bloomsbury. Série se stala celosvětovým fenoménem.

\subwection{Další autorova tvorba:}
\noindent Celá série \textsc{Harry Potter} (7 dílů), \textsc{Fantastická zvířata a kde je najít}, detektivní série pod pseudonymem Robert Galbraith (\textsc{Cormoran Strike}).

\subwection{Aktuálnost tématu a zpracování tématu:}
\noindent Témata přátelství, odvahy, boje se zlem a nalezení vlastní identity zůstávají nadčasová. Kniha popularizovala fantasy pro děti a mládež a inspirovala celou generaci čtenářů.

\vspace*{2em}\\
\wection{OSOBNÍ NÁZOR}
\noindent První díl Harryho Pottera je dokonalý start do kouzelného světa – chytne od první stránky. Nejvíce mě baví kontrast mezi šedivým životem u Dursleyových a barevnými Bradavicemi, stejně jako trio Harry–Ron–Hermiona, které funguje skvěle. Kámen mudrců je chytrý MacGuffin a finále s Quirrellem/Voldemortem napínavé. Rowlingová mistrně mísí humor, napětí a dojemné momenty. I po letech je to jedna z nejpřístupnějších a nejkouzelnějších fantasy sérií. Rozhodně doporučuji – otevírá bránu do světa, kde je přátelství a odvaha silnější než temná moc.

\vfill

\noindent\begin{minipage}{\textwidth}
    {\textcolor{\wpagecolor}{\rule{\linewidth}{0.4pt}}
    \footnotesize
    \textbfhl{Fantasy --} žánr s prvky magie, mytologie a dobrodružství; 
    \textbfhl{Coming-of-age --} příběh dospívání a hledání identity; 
    \textbfhl{MacGuffin --} důležitý předmět, který pohání děj (Kámen mudrců).
    }
\end{minipage}

\newpage

\end{document}